\subsection{仿射 Springer 表示}
本小节引入仿射情形下 Springer 的 $W$-作用在 $\cohog{*}{\sB_{e}}$ 上的类比物。

\subsubsection{仿射 Weyl 群} 把 $W$ 视为余权格 $\xcoch(T)$ 的自同构群，其中 $T$ 是 $G$ 的固定极大环面。{\em 扩张仿射 Weyl 群} $\tilW$ 是半直积
\begin{equation*}
\tilW=\xcoch(T)\rtimes W.
\end{equation*}
当 $G$ 单连通时，$\xcoch(T)$ 由余根张成，此时 $\tilW$ 是 Coxeter 群，其单反射 $\{s_{0},s_{1}, \cdots,s_{r}\}$ 与 $G$ 的扩张 Dynkin 图的节点一一对应。一般情形下，$\tilW$ 是{\em 仿射 Weyl 群} $\Wa=(\ZZ R^{\vee})\rtimes W$（其中 $\ZZ R^{\vee}$ 是余根格）与交换群 $\Om\cong \xcoch(T)/\ZZ R^{\vee}$ 的半直积；$\Wa$ 是 Coxeter 群。$\tilW$ 自然地通过仿射变换作用在仿射空间 $\xcoch(T)_{\RR}$ 上，其中 $\xcoch(T)$ 通过平移作用。

\begin{theorem}[Lusztig {\cite{L96}}，Sage {\cite{Sage}}]\label{th:aff action} 存在 $\tilW$ 在 $\homog{*}{\sY_{\gamma}}$ 上的典范作用。
\end{theorem}

由于 $\sY_{\gamma}$ 不是有限型的，$\ell$-进同调 $\homog{*}{\sY_{\gamma}}$ 应理解为归纳极限 $\varinjlim_{n}\homog{*}{\sY_{\gamma,n}}$，其中把 $\sY_{\gamma}$ 表为射影子概形之并 $\sY_{\gamma,n}$。

\subsubsection{$\tilW$-作用构造梗概} 只考虑 $G$ 单连通的情形，此时 $\tilW=\Wa$ 由仿射单反射 $s_{0}, \cdots, s_{r}$ 生成。对每个 parahoric 子群 $\bP\subset LG$，都有相应的仿射 Springer 纤维 $\tX_{\bP,\gamma}$。当 $\bP$ 含固定的 Iwahori 子群 $\bI$ 时，有投影 $\pi_{\bP,\gamma}:\tY_{\gamma}\to \tX_{\bP,\gamma}$。

设 $L_{\bP}$ 是 $\bP$ 的 Levi 商，$\frl_{\bP}=\Lie L_{\bP}$。有赋值映射 $\ev_{\bP,\gamma}: \tX_{\bP,\gamma}\to [\frl_{\bP}/L_{\bP}]$，定义如下：对满足 $\Ad(g^{-1})\gamma\in\Lie \bP$ 的 $[g]\in\Fl_{\bP}$，把陪集 $[g]=g\bP$ 送到 $\Ad(g^{-1})\gamma$ 在投影 $\Lie \bP\to\frl_{\bP}$ 下的像。此定义在 $L_{\bP}$ 的伴随作用下不变。有 Cartesian 方块
\begin{equation*}
\xymatrix{\tY_{\gamma}\ar[d]^{\pi_{\bP,\gamma}}\ar[r]^{\ev_{\bI,\gamma}} & [\wt{\frl}_{\bP}/L_{\bP}]\ar[d]^{\pi_{\frl_{\bP}}}\\
\tX_{\bP,\gamma}\ar[r]^{\ev_{\bP,\gamma}} & [\frl_{\bP}/L_{\bP}]}
\end{equation*}
其中 $\pi_{\frl_{\bP}}$ 是约化群 $\frl_{\bP}$ 的 Grothendieck 修改。由 $L_{\bP}$ 的 Springer 理论，$W(L_{\bP})$ 作用在直像复形 $\bR\pi_{\frl_{\bP,*}}\DD$ 上（其中 $\DD$ 表示 $[\wt{\frl}_{\bP}/L_{\bP}]$ 的对偶化复形）。由真基变换，得到 $W(L_{\bP})$ 在 $\bR\pi_{\bP,\gamma,*}\DD_{\tY_{\gamma}}$ 上的作用，从而得到它在 $\homog{*}{\sY_{\gamma}}$ 上的作用。取扩张 Dynkin 图第 $i$ 个节点对应的标准 parahoric 子群 $\bP$，此时 $W(L_{\bP})=\jiao{s_{i}}$，于是得到作用在 $\homog{*}{\sY_{\gamma}}$ 上的对合 $s_{i}$。为验证相邻节点 $i$ 与 $j$ 对应的 $s_{i}$ 与 $s_{j}$ 之间的辫关系，可以取标准 parahoric 子群 $\bP$ 使 $W(L_{\bP})=\jiao{s_{i},s_{j}}$，此时辫关系成立，因为 $W(L_{\bP})$ 作用在 $\homog{*}{\sY_{\gamma}}$ 上。这就说明 $\tilW$ 作用在 $\homog{*}{\sY_{\gamma}}$ 上。\qed

尽管 $\tilW$-作用的构造很简单，这些作用的计算却相当困难。这里的一个新现象是：$\tilW$ 在 $\homog{*}{\sY_{\gamma}}$ 上的作用{\em 未必半单}，下面的例子即可看到这一点。

\subsubsection{$\SL_{2}$ 中的一个例子}\label{sss:tilW SL2 ram} 考虑 $G=\SL_{2}$ 与元素 $\gamma=\mat{0}{t^{2}}{t}{0}$。这是次正则情形。仿射 Springer 纤维 $\sY_{\gamma}$ 有两个不可约分支 $C_{0}$ 与 $C_{1}$，两者都同构于 $\PP^{1}$。这里 $C_{0}$ 参数化格链 $\{\L_{i}\}$，其中 $\L_{0}$ 是标准格 $\calO^{2}_{F}$，而 $\L_{-1}\subset \L_{0}$ 变动。另一分支 $C_{1}$ 参数化格链 $\{\L_{i}\}$，其中 $\L_{-1}=t\calO_{F}\oplus\calO_{F}$，而 $\L_{0}$ 变动。基本类 $[C_{0}], [C_{1}]$ 给出最高同调群 $\homog{2}{\sY_{\gamma}}$ 的自然基。可以证明，仿射 Weyl 群 $\tilW=\jiao{s_{0},s_{1}}$ 在 $\homog{2}{\sY_{\gamma}}$ 上的作用在基 $[C_{0}]$ 与 $[C_{1}]$ 下具有如下形式：
\begin{equation*}
s_{0}=\mat{-1}{2}{0}{1}; \quad s_{1}=\mat{1}{0}{2}{-1}.
\end{equation*}
由此可见，$\homog{2}{\sY_{\gamma}}$ 是 $\tilW$ 的符号表示被平凡表示（由 $[C_{0}]+[C_{1}]$ 张成）的{\em 非平凡扩张}。还可以典范地把 $\tilW$-模 $\homog{2}{\sY_{\gamma},\ZZ}$ 等同于环路群 $LG$ 的仿射余根格。

\subsection{评注与推广}

\subsubsection{与轨道积分的关系} 如 \S\ref{s:orb} 将看到的，仿射 Springer 纤维的上同调与点数，同 $p$-进群 $G(\FF_{q}\lr{t})$ 上的轨道积分密切相关。

\subsubsection{扩张对称性} $\tilW$ 在 $\homog{*}{\sY_{\gamma}}$ 上的作用可以扩张为圈积 $\Sym(\xch(T))\rtimes \tilW$ 的作用。对{\em 齐性}仿射 Springer 纤维 $\sY_{\gamma}$（即带有来自圈旋转的环面作用的那些），等变上同调群 $\upH^{*}_{\Gm}(\sY_{\gamma})$ 带有{\em 分级双仿射 Hecke 代数}的作用，后者是 $\Sym(\xch(T))\rtimes \tilW$ 的形变。细节参见 \cite{OY}。Vasserot 与 Varagnolo \cite{Vass} \cite{VV} 在仿射 Springer 纤维的 $K$-群上构造了{\em 双仿射 Hecke 代数}的作用。

\subsubsection{群版本} 把 $\gamma$ 取在 $G(F)$ 中而非 $\frg(F)$ 中，同样可以定义仿射 Springer 纤维的群版本，仍记为 $\tX_{\gamma}$，其约化结构记为 $\sX_{\gamma}$。对 $k$ 代数 $R$，有
\begin{equation}\label{gp AFS}
\tX_{\gamma}(R)=\{[g]\in \Gr_{G}(R)|g^{-1}\gamma g\in L^{+}G(R)\}.
\end{equation}
不过在群版本中，上述定义有一个有趣的推广。回忆 $LG$ 中的 $L^{+}G$-双陪集以支配余权 $\l\in\xcoch(T)^{+}$ 为指标。对 $\l\in\xcoch(T)^{+}$，把相应的双陪集记为 $(LG)_{\l}$，它是 Schubert 分层 $\Gr_{G,\l}$ 在投影 $LG\to \Gr_{G}$ 下的原像。类似地，可以定义 $(LG)_{\leq \l}$ 为 $\Gr_{G,\l}$ 的闭包 $\Gr_{G,\leq \l}$ 的原像。把 \eqref{gp AFS} 中的条件 $g^{-1}\gamma g\in L^{+}G(R)$ 换成 $g^{-1}\gamma g\in (LG)_{\l}(R)$ 或 $g^{-1}\gamma g\in (LG)_{\leq\l}(R)$，并取约化结构，就得到约化的广义仿射 Springer 纤维 $\sX_{\l,\gamma}$ 与 $\sX_{\leq\l, \gamma}$。有开嵌入 $\sX_{\l,\gamma}\incl\sX_{\leq\l,\gamma}$，其补集是支配余权 $\mu\leq\l$ 对应的 $\sX_{\mu,\gamma}$ 之并。引入 $\sX_{\leq\l,\gamma}$ 的动机是给 $G(F)$ 上球 Hecke 函数的轨道积分以几何解释。

A. Bouthier 建立了 $\sX_{\l,\gamma}$ 的基本几何性质，与定理 \ref{th:ft quot} 和定理 \ref{th:aff dim} 平行。

\begin{theorem}[Bouthier \cite{Bouthier}] 设 $\gamma\in G(F)$ 正则半单，$\l\in\xcoch(T)^{+}$。
\begin{enumerate}
\item 广义仿射 Springer 纤维 $\sX_{\l,\gamma}$ 非空，当且仅当 $[\nu_{\gamma}]\leq\l$，其中 $[\nu_{\gamma}]\in\xcoch(T)^{+}_{\QQ}$ 是 $\gamma$ 的 Newton 点，见 \cite[\S2]{KV}。
\item 归纳概形 $\sX_{\l,\gamma}$ 局部有限型。
\item 有
\begin{equation*}
\dim\sX_{\l,\gamma}=\jiao{\rho,\l}+\frac{1}{2}(\val_{F}\Delta(\gamma)-c(\gamma)),
\end{equation*}
其中 $\rho\in\xch(T)_{\QQ}$ 是正根之和的一半，$\Delta(\gamma)$ 与 $c(\gamma)$ 的定义与 Lie 代数情形类似。
\end{enumerate}
\end{theorem}
此定理的证明用到 Vinberg 半群理论，后者是 $G$ 的一种紧化。

\subsubsection{} 在 \cite{KV} 中，Kottwitz 与 Viehmann 对 Lie 代数 $\frg(F)$ 中的元素 $\gamma$ 定义了仿射 Springer 纤维的两种推广。

\subsubsection{} 作为 Hessenberg 簇的类比，还可以考虑如下情形。设 $(\rho,V)$ 是 $k$ 上约化群 $G$ 的线性表示。设 $\L\subset V\otimes F$ 是在 $G(\calO_{F})$ 下稳定的 $\calO_{F}$-格。对 $v\in V\otimes F$，可以定义 $\Gr_{G}$ 的子归纳概形 $\tX_{v}$：
\begin{equation*}
\tX_{\L, v}(R)=\{[g]\in\Gr_{G}(R)|\rho(g^{-1})v\in\L\widehat{\otimes}_{k}R\}.
\end{equation*}
设 $\sX_{\L,v}$ 为 $\tX_{\L,v}$ 的约化结构。这些归纳概形的上同调与相对迹公式中出现的轨道积分相关。

\subsection{习题}

\subsubsection{}\label{exII:SO parahoric} 设 $G=\SO(V,q)$，其中 $V$ 是 $k$ 上配有二次形式 $q$ 的向量空间。仿照 \S\ref{sss:Sp parahoric} 的风格，用 $V\otimes_{k}F$ 中的自对偶格链来解释 $LG$ 的 parahoric 子群与仿射部分旗簇。

\subsubsection{}\label{exII:SL2} 验证 \S\ref{P1 chain} 与 \S\ref{SL2 ram dim 1} 中对 $G=\SL_{2}$ 给出的仿射 Springer 纤维的描述。

\subsubsection{} 设 $G=\SL_{2}$ 与 $\gamma=\mat{0}{t^{n}}{1}{0}\in\frg(F)$。描述 $\sY_{\gamma}$。

\subsubsection{} 设 $G=\SL_{2}$ 与 $\gamma=\mat{0}{t^{n}}{1}{0}\in\frg(F)$。在 $\sX_{\gamma}$ 上构造一个涉及圈旋转（即该作用在 $t$ 上作伸缩）的非平凡 $\Gm$-作用，并确定其不动点。

\subsubsection{} 设 $G=\SL_{2}$ 与 $\gamma=\mat{t^{n}}{0}{0}{-t^{n}}\in\frg(F)$。设 $T\subset G$ 为对角环面，则 $G_{\gamma}=T\otimes_{k}F$。正则轨迹 $\sX^{\reg}_{\gamma}$（见 \S\ref{dim formula discussion}）是什么？研究 $L^+T$ 在 $\sX_{\gamma}$ 上的轨道。

\subsubsection{} 在 \S\ref{sss:L gamma} 的设定下，证明 $\L_{\gamma}$ 在 $\Gr_{G}$ 上的作用自由，从而它在 $\sX_{\gamma}$ 上的作用也自由。再证明 $\L_{\gamma}$ 在 $\sX_{\gamma}$ 的不可约分支集合上的置换作用也自由。

\subsubsection{} 对 $G=\GL_{n}$，设 $L\subset G$ 是由块对角矩阵构成的 Levi 子群，各块大小为 $n_{1},\cdots, n_{s}$，$\sum_{i}n_{i}=n$。设 $\gamma=(\gamma_{1},\cdots,\gamma_{s})\in\frl(F)$ 作为 $\frg(F)$ 中的元素正则半单。不变量 $\Delta^{G}_{L}(\gamma)$（见 \eqref{Delta G L}）用各 $\gamma_{i}$ 的特征多项式的常见不变量如何表示？

\subsubsection{}\label{exII:SL3} 设 $G=\SL_{3}$ 与 $\gamma=\diag(x_{1}t,x_{2}t,x_{3}t)\in\frg(F)$，其中 $x_{i}\in k$ 两两不同且 $x_{1}+x_{2}+x_{3}=0$。描述仿射 Springer 纤维 $\sX_{\gamma}$ 与 $\sY_{\gamma}$。

注：若你有一整天可以消磨，这是一道好习题。

\subsubsection{} 对 $G=\Sp_{6}$ 与例 \ref{ex:Sp} 中的 $\gamma=\mat{0}{tX}{Y}{0}$，描述 $\sX_{\gamma}$ 与 $\sY_{\gamma}$ 上的 $\Gm$-不动点。

\subsubsection{} 验证 \S\ref{sss:tilW SL2 ram} 中的计算。

\subsubsection{} 设 $G=\SL_{2}$ 与 $\gamma=\mat{t}{0}{0}{-t}$。描述仿射 Springer 纤维 $\sY_{\gamma}$。在以 $\sY_{\gamma}$ 的不可约分支给出的基下，$\tilW=\jiao{s_{0},s_{1}}$ 在 $\homog{2}{\sY_{\gamma}}$ 上的作用是什么？
